\documentclass[11pt]{letter}
\usepackage[margin=2.5cm]{geometry}

\begin{document}

\begin{letter}{The Editors \\ \textit{Journal of Fourier Analysis and Applications}}

\opening{Dear Editors,}

We submit for your consideration our manuscript entitled
\textbf{``Spectral gap functions bounded below by band-limited functions:
a 2D uniqueness phenomenon''}
for publication in the \textit{Journal of Fourier Analysis and Applications}.

We study functions $h$ on $\mathbb{R}^2$ whose Fourier transform is
supported outside the unit square and which satisfy the one-sided bound
$h \geq -g$, where $g$ is a non-negative band-limited function. Using
linear programming, we prove that all bounded linear functionals of $h$
vanish, forcing $h = 0$ in $L^1$. The rate of convergence is
$V(N) = 35.5/N^2$, verified for 20 grid sizes up to $N = 128$
($R^2 = 0.9997$).

The result is genuinely two-dimensional: the analogous 1D problem admits
nontrivial solutions. The mechanism involves cooperation between three
spectral bands that individually violate positivity but together satisfy
it. This complements the work of Lagarias and Rodgers (2019), who showed
that without positivity, band-limited mimicry of point processes is
always possible; our result shows that positivity restores uniqueness.

We believe this work is appropriate for \textit{JFAA} because it
addresses a natural question in the Beurling--Selberg tradition of
one-sided approximation by band-limited functions, with connections to
uncertainty principles and determinantal point processes.

A companion paper develops the application of this result to the
Riemann Hypothesis via the Rudnick--Sarnak theorem.

We confirm that this manuscript has not been published previously and is
not under consideration elsewhere.

\closing{Sincerely,}

David Alarc\'on\\
Universidad Pablo de Olavide, Sevilla, Spain\\
\texttt{dalarub@upo.es}

\end{letter}
\end{document}
